\section{Related Work}

\paragraph{Artificial life systems.}
Artificial life systems often study digital organisms, population dynamics,
lineage, self-organization, and open-endedness. SAEE aligns with this
coordinate by treating population evolution, lineage, attractors, and bounded
regime behavior as primary objects of study. Its distinction is that the
paper-facing object explicitly includes mutable operator interpretation,
reflexive coupling, and identity constraints. It does not claim to prove
open-ended evolution.

\paragraph{Evolutionary computation.}
Evolutionary computation studies population search using mutation, selection,
fitness, and recombination operators. SAEE is not framed as a superior
optimizer. It uses evolutionary-computation vocabulary as part of a scientific
object description: operators and selection topology are interpreted as parts
of the system state rather than only as task-level design tools.

\paragraph{Complex systems.}
Complex systems research provides the language of attractors, regimes,
stability, phase spaces, and transitions. SAEE's current frozen phase-space
fits this framing: the observed dynamics are dominated by a stable lineage
basin, with no observed collapse or cross-regime transition under the current
constraints.

\paragraph{Dynamical systems.}
Dynamical systems theory supplies the state-transition framing for the formal
surface. SAEE contributes a local reflexive evolutionary object in which
population state, transformation, selection, lineage, observer state, and
identity constraint are represented together. This is not a claim of a new
universal dynamical systems theory or externally validated physical law.
